% Lecture notes / outline for RankAlign talk
% Edit this, then feed it back to get reveal.js slides
%
% TIP: Add visualization cues like [VIZ: description] where you want
% interactive elements. The more specific, the better the result.
%
% STYLE GUIDE for reveal.js slides:
%   - Warm cream/beige background (#f5f0e8), dark text (#2d2d2d)
%   - Clean sans-serif font (Source Sans 3 or similar from Google Fonts)
%   - Headings: bold, dark, no all-caps transform
%   - Equation boxes: white background, subtle shadow, left accent border (warm rose)
%   - Color-code equation terms where helpful (blue for one quantity, orange for another, gray for labels)
%   - Keep equations compact — use \small or \scriptstyle to avoid overflow
%   - IMPORTANT: Always check that content fits within 960×700px slide area.
%     Stack elements vertically rather than side-by-side if equations are wide.
%     Long fractions/logs especially tend to overflow — prefer stacked layout for these.
%   - Cards/boxes: white with slight shadow, rounded corners
%   - Interactive elements: white card backgrounds, muted colored accents
%   - Buttons: soft rounded, muted fills, dark text
%   - Canvas charts: dark axes/text on light background (no white-on-dark)
%   - Overall vibe: academic but modern, like a polished course slide deck
%   - Background: white (#ffffff)
%
% REFERENCES:
%   - Put citations at the bottom of the slide in small gray font
%   - Use class="citation" for the container (positioned at bottom-left)
%   - Format: "Author et al., Short Title. Venue Year"
%   - Only cite on slides where the work is first introduced or directly relevant
%
% SLIDE NUMBERING (for reference when editing):
%   Slide 0: Title
%   Slides 1–3: Section 1 — The Generator–Validator Gap
%     Slide 1: Two probing modes (generator/validator prompt boxes)
%     Slide 2: Formal definitions (gen/val equations)
%     Slide 3: Interactive scatter plot (category buttons)
%   Slide 4: The Generator–Validator Gap (V2G + G2V losses)
%   Slide 5: Current Direction (combined loss, research questions)
%   Slide 6: Tasks and Datasets (Hypernymy, PlausibleQA, AmbigQA, IFEval)
%   Slide 7: Frequency Effect (car example)
%   Slides 8–9: Section 2 — The Typicality Confound
%     Slide 8: Typicality decomposition + PMI
%     Slide 9: Interactive slider for b

\documentclass[11pt]{article}
\usepackage[margin=1in]{geometry}
\usepackage{amsmath,amssymb}

\title{The Generator-Validator Gap and Frequency Correction}
\author{Juan Diego Rodriguez}
\date{}

\begin{document}
\maketitle

% =====================================================================
\section{The Generator--Validator Gap}
% =====================================================================

Language models can be probed for the same knowledge in two ways.
Example: hypernymy (``a bee is a kind of \underline{\quad}'').

\begin{itemize}
  \item \textbf{Generator:} ``A bee is a kind of'' $\to$ look at next-token probabilities.
  \item \textbf{Validator:} ``Is a bee a kind of furniture? Answer:'' $\to$ look at Yes/No log-odds.
\end{itemize}

A well-calibrated model should rank these consistently --- but often it doesn't.

Formally:
\begin{align}
  \text{gen}(y|x) &= \log P_{\text{LLM}}(y \mid x) \\
  \text{val}(y|x) &= \log \frac{P(\text{Yes} \mid \text{``Is } y \text{ correct?''})}{P(\text{No} \mid \cdot)}
\end{align}

We measure the gap via Pearson correlation between gen and val scores across candidate answers for each question.

% [VIZ: Scatter plot of gen vs val scores for a hypernymy example,
%  color-coded by correctness. Allow clicking between categories
%  (insect, furniture, vehicle) to see how the gap varies.]


% =====================================================================
\section{The Typicality Confound}
% =====================================================================

Raw generator scores are confounded by unconditional word frequency (``typicality''):
\[
  \log P(y \mid x) \;\approx\; \underbrace{\log P_{\text{val}}(y|x)}_{\text{correctness}} + \underbrace{\log P(y)}_{\text{typicality}}
\]

Common words get high generator scores regardless of correctness.

Standard fix: Pointwise Mutual Information (PMI):
\[
  \text{PMI}(y; x) = \log P(y|x) - \log P(y)
\]

We estimate typicality using GPT-2: $\text{typ}(y) = \log P_{\text{GPT2}}(y)$.

More generally, we can consider $\text{gen}(y|x) - b \cdot \text{typ}(y)$ for various $b$.

% [VIZ: Slider for b in gen - b*typ. Show how correlation and GenROC
%  change as b varies from 0 to 3, with actual data points at b=0,0.5,1,1.5,2,3.
%  Mark b=1.5 as the optimum.]


% =====================================================================
\section{Fitted Corrections Overfit}
% =====================================================================

With only $\sim$14 candidates per question, fitting $b$ per-prompt is catastrophic.

Key result: strict monotonic relationship between parameter count and LOO performance:
\begin{center}
\begin{tabular}{cc}
  \# Parameters & Best LOO GenROC \\
  \hline
  0 & 76.0\% (PMI$_{1.5}$) \\
  1 & 66.3\% \\
  2 & 63.4\% \\
  4 & 60.8\% \\
  5 & 56.4\% \\
\end{tabular}
\end{center}

The in-sample vs LOO gap is dramatic (e.g.\ 73.4\% in-sample $\to$ 20.1\% LOO for a 3-param method).

Takeaway: $b \approx 1.5$ with zero parameters is the best correction.

% [VIZ: Bar chart comparing in-sample vs LOO for each method,
%  with drop annotations. Or: scatter of #params vs LOO GenROC
%  with a "sweet spot" region around 0 params.]


% =====================================================================
\section{RankAlign}
% =====================================================================

Train the model so its validator rankings align with a reference signal.

Pipeline:
\begin{enumerate}
  \item Precompute gen \& val scores for all candidates.
  \item Build pairs from the same question where score gap exceeds a threshold $\delta$.
  \item Pairwise logistic loss: $-\log \sigma(s_j - s_i - \Delta_{\text{ref}})$.
  \item LoRA fine-tuning.
\end{enumerate}

Three loss variants:
\begin{itemize}
  \item \textbf{Pref:} preference loss only.
  \item \textbf{Comb:} preference + NLL$_{\text{val}}$ + NLL$_{\text{gen}}$.
  \item \textbf{SFT:} NLL only, no preference.
\end{itemize}

Training modes:
\begin{itemize}
  \item d2g: align validator $\to$ generator signal.
  \item g2d: align generator $\to$ validator signal.
\end{itemize}

% [VIZ: Pipeline flow diagram showing the four steps.
%  Maybe animated, highlighting each step in sequence.]


% =====================================================================
\section{Semi-Supervised Extension}
% =====================================================================

What if only 10\% of prompts have ground-truth labels?

\begin{itemize}
  \item \textbf{Label-only:} Discard unlabeled prompts $\to$ smaller training set.
  \item \textbf{Semi-supervised:} Labeled prompts get full loss; unlabeled prompts get preference loss only (NLL terms zeroed).
\end{itemize}

Key insight: unlabeled prompts still carry useful \emph{preference} signal --- the model can learn relative rankings even without ground truth.

Split is at the prompt level (all answers for a question are either labeled or unlabeled). Deterministic via seed for reproducibility.

% [VIZ: Grid of training prompts, slider to adjust labeled ratio.
%  Show labeled (green check) vs unlabeled (?) prompts.
%  Side-by-side boxes explaining label-only vs semi-supervised.]


% =====================================================================
\section{Results}
% =====================================================================

Across tasks (PlausibleQA, AmbigQA, Hypernym ID/OOD), semi-supervised consistently outperforms label-only at 10\% labels.

Selected results (Pearson $r$, typicality-corrected):
\begin{center}
\begin{tabular}{lcc}
  & Comb+v label-only & Comb+v semi \\
  \hline
  PlausibleQA & 0.359 & 0.393 \\
  AmbigQA & 0.249 & 0.380 \\
  Hypernym (ID) & 0.533 & 0.763 \\
  Hypernym (OOD) & 0.491 & 0.709 \\
\end{tabular}
\end{center}

The effect is largest on Hypernym, where semi doubles the correlation.

% [VIZ: Bar chart per task, switchable between tasks.
%  Show label-only vs semi bars with base model as dashed line.
%  Use actual numbers from semisupervised.tex tables.]


% =====================================================================
\section{Key Findings}
% =====================================================================

\begin{enumerate}
  \item Semi-supervised consistently beats label-only (up to +0.50 correlation on Hypernym).
  \item PMI$_{1.5}$ is the best parameter-free typicality correction (53.5\% corr, 76.0\% GenROC).
  \item Fitted corrections overfit: more parameters $\to$ worse generalization.
  \item Length is largely redundant with typicality (within-prompt $R^2 = 0.79$).
\end{enumerate}


% =====================================================================
\section{Datasets}
% =====================================================================

\begin{itemize}
  \item \textbf{Hypernymy:} ``A \_\_ is a kind of \_\_'' --- 8 in-domain + 10 OOD nouns.
  \item \textbf{PlausibleQA:} Open-domain questions with plausible distractors, GPT-4 labels. 33 questions $\times$ $\sim$14 candidates.
  \item \textbf{AmbigQA:} Ambiguous questions with multiple valid answers. GPT-4o generated negatives.
  \item \textbf{IFEval:} Instruction-following evaluation, long-form completions.
\end{itemize}

Model: google/gemma-2-2b with LoRA.


% =====================================================================
\section{Summary}
% =====================================================================

\begin{itemize}
  \item The G-V gap is real and measurable.
  \item Typicality is the main confound; PMI$_{1.5}$ is the best simple fix.
  \item RankAlign closes the gap via ranking-based training.
  \item Semi-supervised training works: 10\% labeled + 90\% unlabeled preference signal matches or exceeds fully supervised.
\end{itemize}

\end{document}
